\scanpage{3}
% archive foliation: 2
\diagram{\begin{tikzcd}[column sep=small]
0\arrow[r]&\Omega\arrow[r]&E\arrow[r]&\mathcal{O}_S\arrow[r]&0
\end{tikzcd}}

en ce sens que le \inserted{\unsure{contre}} foncteur qui à une telle extension fait correspondre $(P,\sigma)$ avec $P=\mathbb{P}(E)$ et $\sigma$ la section définie par le quotient \inserted{\unsure{inv.}} $\mathcal{O}_S$ de $E$, est une équivalence de catégories.

On trouve que les $\mathcal{O}_P/\mathcal{J}^{n+1}$ (où $\mathcal{J}$ est l'idéal de $\sigma(S)$ dans $P$) s'identifient comme $\mathcal{O}_S$-algèbres augmentées à
\[
\operatorname{Sym}^n(
\]
